\par\medskip
\begin{center}\hypertarget{native-163-article116}{}\pdfbookmark[2]{116.}{article-116}116.\end{center}
Antequam determinatio alicuius orbitae ex observationibus geocentricis suscipitur, his quaedam reductiones applicandae erunt, propter nutationem, praecessionem, parallaxin et aberrationem, siquidem summa praecisio requiritur: in crassiori enim calculo has minutias negligere licebit.

Planetarum et cometarum observationes vulgo expressae proferuntur per ascensiones rectas et declinationes apparentes, i. e. ad situm aequatoris apparentem relatas. Qui situs quum propter nutationem et praecessionem variabilis adeoque pro diversis observationibus diversus sit, ante omnia loco plani variabilis planum aliquod fixum introducere conveniet, ad quem finem vel aequator situ suo medio pro aliqua epocha, vel ecliptica adoptari poterit: planum posterius plerumque adhiberi solet, sed prius quoque commodis peculiaribus haud spernendis se commendat.

Quoties itaque planum aequatoris eligere placuit, ante omnia observationes a nutatione purgandae, ac dein adhibita praecessione ad epocham quandam arbitrariam reducendae sunt: haec operatio prorsus convenit cum ea, per quam e loco stellae fixae observato eiusdem positio media pro epocha data derivatur, adeoque explicatione hic non indiget. Sin vero planum eclipticae adoptare constitutum est, duplex methodus patebit: scilicet vel ex ascensionibus rectis et declinationibus ob nutationem et praecessionem correctis deduci poterunt longitudines et latitudines adiumento obliquitatis mediae, unde longitudines iam ad aequinoctium medium relatae prodibunt; vel commodius ex ascensionibus rectis et declinationibus apparentibus adiumento obliquitatis apparentis computabuntur longitudines et latitudines, ac dein illae a nutatione et praecessione purgabuntur.

Loci terrae singulis observationibus respondentes per tabulas Solares computantur, manifesto autem ad idem planum referendi erunt, ad quod observationes corporis coelestis relatae sunt. Quamobrem in computo longitudinis Solis negligetur nutatio; dein vero haec longitudo adhibita praecessione ad \SourceLeaf{164}epocham fixam reducetur, atque 180 gradibus augebitur; latitudini Solis, siquidem eius rationem habere operae pretium videtur, signum oppositum tribuetur: sic positio terrae heliocentrica habebitur, quam, si aequator pro plano fundamentali electus est, adiumento obliquitatis mediae in ascensionem rectam et declinationem transformare licebit.
